\chapter{Trabalhos Relacionados}

Este capítulo apresenta e discute ferramentas oferecidas na indústria por SGBDs e também
propostas na literatura acadêmica relacionadas com a importação de dados por SGBDs
heterogêneos. A preferência está em dados no formato CSV que devem ser importados por
alguma tecnologia de banco de dados, que é o foco deste trabalho.

\section{Importação de Arquivo CSV para um Único SGBD}

Esta seção investiga mecanismos de importação de dados presentes em SGBDs populares na
indústria que são considerados neste trabalho.

\subsection{PostgreSQL}

No SGBD PostgreSQL você utiliza o comando \texttt{COPY} para realizar importações de dados
CSV \cite{postgresql_copy}. É necessário especificar a tabela com os nomes das colunas após a cláusula
\texttt{COPY}. A ordem das colunas deve ser a mesma que aquelas no arquivo CSV, conforme
ilustra o exemplo a seguir:

\begin{lstlisting}[style=sqlstyle]
COPY nome_da_tabela(col1, col2, col3)
FROM 'C:/caminho_do_arquivo/dados_de_amostra.csv'
DELIMITER ','
CSV HEADER;
\end{lstlisting}

Caso o arquivo CSV contenha todas as colunas da tabela, você não precisa especificá-las
explicitamente:

\begin{lstlisting}[style=sqlstyle]
COPY nome_da_tabela
FROM 'C:/caminho_do_arquivo/dados_de_amostra.csv'
DELIMITER ','
CSV HEADER;
\end{lstlisting}

Você também deve inserir o caminho do arquivo CSV após a palavra-chave \texttt{FROM}. Como
o formato do arquivo é CSV, você precisa especificar \texttt{DELIMITER}, bem como a
cláusula \texttt{CSV}. A cláusula \texttt{HEADER} deve ser especificada para indicar que o
arquivo CSV contém um cabeçalho. Quando o comando \texttt{COPY} importa dados, ele ignora
o cabeçalho do arquivo.

O arquivo CSV deve ser lido diretamente pelo servidor PostgreSQL, não pela aplicação
cliente. Portanto, ele deve ser acessível pela máquina do servidor PostgreSQL. Além disso,
é necessário ter permissão de superusuário para executar com sucesso o comando
\texttt{COPY}.

\subsection{Redis}

A maneira preferida de importar uma massa de dados para o Redis é gerar um arquivo de
texto contendo o protocolo Redis, em formato bruto, para chamar os comandos necessários
para inserir os dados requeridos \cite{redis_bulkloading}.

Por exemplo, se é necessário gerar um grande conjunto de dados com bilhões de chaves no
formato ``chaveN $\rightarrow$ ValorN'', deve ser criado um arquivo contendo os seguintes
comandos:

\begin{lstlisting}[style=textstyle]
SET Chave0 Valor0
SET Chave1 Valor1
...
SET ChaveN ValorN
\end{lstlisting}

Uma vez criado esse arquivo, a ação restante é alimentá-lo no Redis. No passado, a maneira
de fazer isso era usar o netcat com o seguinte comando:

\begin{lstlisting}[style=bashstyle]
(cat dados.txt; sleep 10) | nc localhost 6379
\end{lstlisting}

Nas versões 2.6 ou posteriores do Redis, o utilitário \texttt{redis-cli} suporta um novo
modo chamado modo \textit{pipe} que foi projetado para realizar o carregamento em massa. O
comando a ser executado neste caso é o seguinte:

\begin{lstlisting}[style=bashstyle]
cat dados.txt | redis-cli --pipe
\end{lstlisting}

O utilitário \texttt{redis-cli} imprime um pequeno resumo. A saída típica é semelhante a
esta:

\begin{lstlisting}[style=textstyle]
All data transferred. Waiting for the last reply...
Last reply received from server.
errors: 0, replies: 1000000
\end{lstlisting}

A linha de contagem ao final é a informação mais importante do resumo: \texttt{errors}
indica quantos comandos foram rejeitados pelo servidor, e \texttt{replies} indica o total
de respostas recebidas. Se \texttt{errors: 0}, todos os comandos foram processados com
sucesso. Um valor maior que zero em \texttt{errors} significa que ao menos um comando
falhou --- seja por sintaxe inválida no arquivo, por tipo de dado incompatível com a chave
existente, ou por qualquer outra razão que o Redis tenha rejeitado.

Vale notar que o \texttt{-{}-pipe} não exibe cada resposta individualmente (como um
\texttt{+OK} por \texttt{SET}), o que é intencional: imprimir bilhões de respostas
destruiria o ganho de \textit{performance} do modo \textit{pipe} e inundaria o terminal.
Em vez disso, o \texttt{redis-cli} lê e contabiliza todas as respostas silenciosamente
conforme chegam, entregando apenas o agregado final.

Por fim, uma verificação útil: o valor de \texttt{replies} deve ser igual ao número de
comandos enviados. Uma divergência entre os dois pode indicar que o arquivo de entrada
estava truncado ou malformado, mesmo que \texttt{errors} apareça como zero.

\subsection{MongoDB}

No SGBD MongoDB, o processo de importação de dados em formato CSV é realizado através do
utilitário \texttt{mongoimport} \cite{mongodb_mongoimport}. Sua sintaxe básica exige que se informe o banco de dados
de destino através do parâmetro \texttt{-{}-db}, seguido do nome da coleção que receberá
os dados, especificado em \texttt{-{}-collection}. Caso esse último parâmetro seja omitido,
o nome da coleção será automaticamente inferido a partir do nome do arquivo CSV fornecido.

\begin{lstlisting}[style=bashstyle]
mongoimport --db nome_do_db --collection nome_da_colecao \
  --type csv --headerline --ignoreBlanks \
  --file caminho/do/arquivo.csv
\end{lstlisting}

O tipo do arquivo de entrada deve ser declarado explicitamente através do parâmetro
\texttt{-{}-type}, que neste caso recebe o valor \texttt{csv}. Para que o
\texttt{mongoimport} reconheça os nomes dos campos a partir do próprio arquivo, utiliza-se
o parâmetro \texttt{-{}-headerline}, que instrui o utilitário a tratar o conteúdo da
primeira linha do CSV como o cabeçalho dos campos. Opcionalmente, o parâmetro
\texttt{-{}-ignoreBlanks} pode ser utilizado para que valores em branco presentes no
arquivo sejam ignorados durante a importação, evitando que campos vazios sejam persistidos
na coleção. Por fim, o caminho do arquivo CSV a ser importado é fornecido através do
parâmetro \texttt{-{}-file}.

\subsection{Cassandra}

A importação de dados CSV para o Cassandra utiliza o comando \texttt{sstableloader} e
envolve vários passos \cite{cassandra_sstableloader}. Inicialmente, é necessário se certificar que o arquivo CSV está no
formato adequado para o Cassandra. Isso geralmente envolve garantir que os tipos de dados
e as colunas correspondem ao esquema da tabela no Cassandra.

O \texttt{sstableloader} requer que os dados estejam em um formato específico chamado
SSTable. Para converter o CSV em SSTables, você pode usar o comando \texttt{cqlsh}
fornecido pelo Cassandra. Um exemplo de uso é o seguinte:

\begin{lstlisting}[style=bashstyle]
cqlsh -e "COPY keyspace_name.table_name TO 'output_directory';"
\end{lstlisting}

Na sequência, é possível usar o \texttt{sstableloader} para carregar os dados. A sintaxe
do comando é a seguinte:

\begin{lstlisting}[style=bashstyle]
sstableloader -d <hostname> -u <username> -pw <password> \
  output_directory/keyspace_name/table_name
\end{lstlisting}

\begin{itemize}
  \item \texttt{<hostname>}: O endereço do nó Cassandra.
  \item \texttt{<username>}: O nome de usuário, se a autenticação estiver habilitada.
  \item \texttt{<password>}: A senha correspondente.
\end{itemize}

Este comando move os dados do diretório de saída para o Cassandra.

\subsection{Neo4j}

Usar \texttt{neo4j-admin} para importar grandes conjuntos de dados no Neo4j envolve alguns
passos \cite{neo4j_import}. Esta ferramenta de linha de comando é projetada para importações em massa
eficientes.

Certifique-se de que seus arquivos CSV estejam estruturados corretamente e sigam o formato
necessário para nós e relacionamentos. Crie um banco de dados não povoado:

\begin{lstlisting}[style=bashstyle]
neo4j-admin create-db --database=nome_do_seu_banco_de_dados \
  --from=diretorio_com_arquivos_csv
\end{lstlisting}

Substitua \texttt{nome\_do\_seu\_banco\_de\_dados} pelo nome desejado para o novo banco de
dados. A opção \texttt{-{}-from} especifica o diretório onde estão localizados seus
arquivos CSV. Execute o comando de importação com \texttt{neo4j-admin}. O comando exato
depende do seu modelo de dados, mas pode se parecer com isto:

\begin{lstlisting}[style=bashstyle]
neo4j-admin import --database=nome_do_seu_banco_de_dados \
  --mode=csv \
  --nodes:Rotulo1 nos.csv \
  --nodes:Rotulo2 nos2.csv \
  --relationships:TIPO_RELACIONAMENTO relacionamentos.csv
\end{lstlisting}

Substitua \textit{Rotulo1}, \textit{Rotulo2}, \textit{nos.csv}, \textit{nos2.csv},
\textit{TIPO\_RELACIONAMENTO} e \textit{relacionamentos.csv} pelos seus rótulos específicos
e nomes de arquivo.

\section{Importação de Arquivo CSV para SGBD Multimodelo}

Ferramentas comerciais e de código aberto tratam de \textbf{migração} ou
\textbf{modelagem} entre tecnologias heterogêneas, embora raramente com o foco deste TCC
(um único CSV largo + regras declarativas + cinco paradigmas).

\begin{itemize}
  \item \textbf{ArangoDB} é um SGBD multimodelo nativo (documento/chave-valor/grafo) com
        \texttt{arangoimport} para CSV; a importação é para \textbf{um} motor multimodelo,
        não para cinco SGBDs independentes.

  \item \textbf{\texttt{mongoimport} / \texttt{COPY} / \texttt{redis-cli -\/-pipe} /
        \texttt{cqlsh} / \texttt{neo4j-admin}} (seção 3.1) cobrem importação \textbf{para
        um único} produto por vez.

  \item \textbf{dbcrossbar} \cite{dbcrossbar2024} copia dados tabulares entre PostgreSQL,
        BigQuery, armazenamento em nuvem e CSV, incluindo conversão de esquemas; a ênfase é
        \textit{pipeline} de migração, não mapeamento semântico de entidades aninhadas com
        filtros por valor.

  \item \textbf{Hackolade} \cite{hackolade2024polyglot} oferece \textit{polyglot data
        modeling} com engenharia de esquema para dezenas de alvos; trata-se de
        \textbf{modelagem visual} e geração de artefatos, e não de um utilitário CLI único
        para importar um CSV operacional com validação prévia de filtros.
\end{itemize}

\section{Considerações sobre os Trabalhos Relacionados}

A literatura sobre \textbf{metamodelos unificados} para SQL e NoSQL (por exemplo, U-Schema
\cite{fernandezcandel2022unified}) e sobre \textbf{evolução de esquema} em ambientes heterogêneos
\cite{chillon2023propagating,chillon2024schema} informa decisões de validação e de representação de entidades
no JSON de configuração, mas não substitui uma ferramenta de \textit{bulk import} orientada
a CSV.

\textit{Surveys} recentes convergem em desafios abertos da gestão poliglota de dados:
migração automatizada entre \textit{stores} quando o \textit{workload} muda, planejamento
de consultas entre sistemas, gestão de esquemas multi-modelo e preservação das propriedades
não funcionais de cada SGBD \cite{kiehn2022polyglot,glake2022towards}.
\citeonline{royhubara2022selecting} tratam da \textbf{seleção de SGBDs} por fragmento de
aplicação; \citeonline{silva2024modelagem} mapeiam o estado da arte em \textbf{modelagem
poliglota}. Nenhum desses trabalhos implementa, contudo, um utilitário declarativo que
receba um CSV operacional único e o materialize simultaneamente em cinco paradigmas
distintos com validação estática prévia.

Abordagens recentes reforçam essa tendência. \citeonline{pereira2025transparent} propõem um
\textit{framework} de persistência poliglota \textbf{transparente}, com uma {API} unificada e
extensível a diferentes tipos de banco, enquanto \citeonline{sachdeva2024plugging} exploram a
combinação de bancos \textbf{multimodelo} e persistência poliglota para acomodar a variedade dos
dados. Ambos os trabalhos atuam, porém, na \textbf{camada de acesso da aplicação} --- oferecendo
uma interface única para ler e gravar em múltiplos SGBDs em tempo de execução ---, e não na
\textbf{carga inicial declarativa} a partir de um CSV operacional, que é o foco da
\textit{Polyglot Import CSV}.

Sistemas \textit{polystore} como o BigDAWG \cite{duggan2015bigdawg} e \textit{middlewares}
de portabilidade em nuvem \cite{kaur2020middleware} tratam de \textbf{consultas federadas} e
movimentação de dados em ecossistemas complexos. \citeonline{tan2017query} classificam
essas soluções (federado, multistore, polystore, poliglota) e avaliam transparência,
autonomia e heterogeneidade --- eixos nos quais importadores nativos (seção 3.1) e
ferramentas de migração (seção 3.2) também divergem, mas sem convergir para um
\textit{pipeline} único de ingestão.

Ferramentas como \textbf{dbcrossbar} \cite{dbcrossbar2024} e \textbf{Hackolade}
\cite{hackolade2024polyglot} aproximam-se do domínio de integração de dados, porém com
foco em \textit{pipelines} de migração ou modelagem visual, respectivamente --- não em
regras declarativas sobre entidades aninhadas, filtros por valor (\texttt{each}) e
validação por JSON Schema antes de qualquer conexão com SGBDs.

\textbf{Lacuna identificada:} a combinação de (i) um CSV largo como fonte operacional,
(ii) configuração JSON validada estaticamente, (iii) materialização distinta por paradigma
(relacional, chave-valor, documento, colunar, grafo) e (iv) modo \textit{dry-run} para
revisão prévia. O \textit{Polyglot Import CSV} ocupa essa lacuna como ferramenta de
\textbf{bootstrap de persistência poliglota} --- carga inicial e reprodutível de dados ---
enquanto \textit{middlewares} e \textit{polystores} assumem dados já residentes ou
consultas em tempo de execução. Trabalhos futuros podem integrar a ferramenta a mecanismos
de migração adaptativa descritos por \citeonline{kiehn2022polyglot} e
\citeonline{glake2022towards}.
